\documentclass[11pt]{article}

\usepackage[margin=1in]{geometry}
\usepackage{amsmath, amssymb}
\usepackage{xcolor}
\usepackage{listings}
\usepackage{hyperref}
\usepackage{array}
\usepackage{booktabs}

\hypersetup{
    colorlinks=true,
    linkcolor=blue!60!black,
    urlcolor=blue!60!black,
    citecolor=blue!60!black
}

\lstdefinestyle{bugpython}{
  language=Python,
  basicstyle=\ttfamily\small,
  keywordstyle=\color{blue!60!black},
  stringstyle=\color{green!40!black},
  commentstyle=\color{gray!70!black},
  showstringspaces=false,
  breaklines=true,
  frame=single,
  rulecolor=\color{gray!30},
  columns=fullflexible,
  keepspaces=true
}

\title{SymPy Bug Report}
\author{}
\date{}

\begin{document}

\maketitle

\section{Bug 6: Piecewise Left-Hand Limit Returns Right Branch}

\subsection{Minimal Reproducer}

\medskip

\begin{lstlisting}[style=bugpython]
from sympy import Piecewise, Rational, limit, symbols

x = symbols("x", real=True)
p = Piecewise((0, x < 0), (1, True))

print(limit(p, x, 0, dir="-"))
print(p.subs(x, Rational(-1, 10)))
print(p.subs(x, Rational(1, 10)))
\end{lstlisting}

\subsection{Observed Behavior}

\medskip

\begin{lstlisting}[style=bugpython]
limit(Piecewise((0, x < 0), (1, True)), x, 0, dir='-') = 1
left sample p.subs(x, -1/10) = 0
right sample p.subs(x, 1/10) = 1
\end{lstlisting}

SymPy 1.14.0 returns \(1\), the value of the default branch, for the left-hand
limit.  The same expression evaluates to \(0\) at a nearby point on the left
side, so the returned value is the value from the wrong side of the jump.

\subsection{Expected Behavior}

The left-hand limit should be \(0\):
\[
  \lim_{x\to 0^-}
  \operatorname{Piecewise}\bigl((0,x<0),(1,\mathrm{True})\bigr)=0.
\]
Software-wise, \texttt{limit(p, x, 0, dir="-")} should select the branch that is
active in a punctured left neighborhood of \(0\), not the branch that is active
at or to the right of the boundary.

\subsection{Mathematical Explanation}

For every real \(x\) with \(x<0\), the condition in the first branch is true, so
\[
  \operatorname{Piecewise}\bigl((0,x<0),(1,\mathrm{True})\bigr)=0.
\]
Therefore, for every sufficiently small negative \(x\), the expression is
identically \(0\).  The one-sided limit from the left is consequently
\[
  \lim_{x\to 0^-} p(x)=0.
\]

The returned value \(1\) is not an equivalent representation or a harmless
branch choice.  It is the right/default branch value: for instance
\[
  p(-1/10)=0,\qquad p(1/10)=1.
\]
A left-hand limit at this jump must be determined by the left-neighborhood
values, not by substituting the boundary point into the branch condition.

\subsection{Source-Level Diagnosis}

The diagnosis identifies the source file
\nolinkurl{sympy/functions/elementary/piecewise.py}.  The relevant method and
traced call path are:

\begin{lstlisting}[style=bugpython]
limit(p, x, 0, dir="-")
Limit(...).doit(deep=False)
Piecewise._eval_as_leading_term(x, logx, cdir)
\end{lstlisting}

The surrounding limit code in \nolinkurl{sympy/series/limits.py} constructs a
\texttt{Limit} object and evaluates it via \texttt{doit}.  For a finite
directed limit, the limit engine shifts the point of approach to \(0\) and
records the left-hand direction as \texttt{cdir = -1}.  It then asks the
\texttt{Piecewise} expression for its leading term.

The relevant method receives the direction parameter but ignores it:

\begin{lstlisting}[style=bugpython]
def _eval_as_leading_term(self, x, logx, cdir):
    for e, c in self.args:
        if c == True or c.subs(x, 0) == True:
            return e.as_leading_term(x)
\end{lstlisting}

For the condition \(x<0\), direct substitution at the boundary gives
\((x<0)\rvert_{x=0}\), which is false.  The method therefore skips the left
branch even though \texttt{cdir = -1} means the approach is through negative
values of the shifted variable.  It then reaches the default \texttt{True}
branch and returns the leading term of the expression \(1\), which the limit
engine reports as the left-hand limit.

A plausible fix direction supported by the diagnosis is to make this
\texttt{Piecewise} leading-term method direction-aware.  For relational
conditions involving the limit variable, it should test which branch holds in a
punctured neighborhood determined by \texttt{cdir}, for example through a
signed dummy substitution or an interval/set containment check, rather than by
substituting the exact boundary point.

\subsection{Additional Failing Instances}

The independent verification checks the representative case together with five
more constant jumps of the same form:
\[
  p_{a,b}(x)=\operatorname{Piecewise}\bigl((a,x<0),(b,\mathrm{True})\bigr).
\]
For each pair \((a,b)\), every sufficiently small negative \(x\) satisfies
\(x<0\), so the expected left-hand limit is \(a\).  SymPy instead returns \(b\),
the default branch value.

\begin{center}
\footnotesize
\renewcommand{\arraystretch}{1.3}
\begin{tabular}{@{}>{\raggedright\arraybackslash}p{0.34\linewidth}>{\raggedright\arraybackslash}p{0.28\linewidth}>{\raggedright\arraybackslash}p{0.28\linewidth}@{}}
\toprule
\textbf{Input / Expression} & \textbf{SymPy behavior} & \textbf{Expected behavior} \\
\midrule
\(p_{0,1}\) & returns \(1\), the default branch & returns \(0\), the left branch \\
\(p_{2,5}\) & returns \(5\), the default branch & returns \(2\), the left branch \\
\(p_{-3,7}\) & returns \(7\), the default branch & returns \(-3\), the left branch \\
\(p_{4,-2}\) & returns \(-2\), the default branch & returns \(4\), the left branch \\
\(p_{11,13}\) & returns \(13\), the default branch & returns \(11\), the left branch \\
\(p_{-8,-1}\) & returns \(-1\), the default branch & returns \(-8\), the left branch \\
\bottomrule
\end{tabular}
\end{center}

The expected behavior is the same across the family because the condition
\(x<0\) is true throughout a punctured left neighborhood of the limit point.
The right/default branch may determine values at \(x=0\) or for \(x>0\), but it
does not determine the left-hand limit.

\subsection{Related Issues Assessment}

The bug-finding harness initially classified this as a new failure mode with no
public precedent, but that classification was incorrect.  This is a known,
long-standing problem in SymPy.  It is the same root cause described in the open
issue \href{https://github.com/sympy/sympy/issues/23836}{\#23836} (``Incorrect
results for limits of Piecewise at discontinuity''), which states that the value
of the \texttt{Piecewise} at the boundary should not be considered by
\texttt{limit}, and it is a sibling of the open issue
\href{https://github.com/sympy/sympy/issues/27236}{\#27236} (``bug with limit()
function''), where one-sided limits of a discontinuous \texttt{Piecewise} return
the wrong branch.  There are also open pull requests targeting the precise
function diagnosed above (\texttt{Piecewise.\_eval\_as\_leading\_term} ignoring
the direction \texttt{cdir}):
\href{https://github.com/sympy/sympy/pull/29589}{\#29589},
\href{https://github.com/sympy/sympy/pull/29538}{\#29538}, and
\href{https://github.com/sympy/sympy/pull/29920}{\#29920}.  The closest related
public material is therefore not merely subsystem-level: these issues and PRs
describe this exact failure.  The bug remains a clear, reproducible instance of
this known family, and the focused regression test below guards branch selection
in directed limits of \texttt{Piecewise} expressions.

\subsection{Proposed SymPy Regression Test}

\medskip

\begin{lstlisting}[style=bugpython]
import pytest

from sympy import Piecewise, limit, symbols


@pytest.mark.parametrize(
    "left_value, default_value",
    [(0, 1), (2, 5), (-3, 7), (4, -2), (11, 13), (-8, -1)],
)
def test_piecewise_left_hand_limit_uses_left_branch(left_value, default_value):
    x = symbols("x", real=True)
    p = Piecewise((left_value, x < 0), (default_value, True))
    assert limit(p, x, 0, dir="-") == left_value
\end{lstlisting}

\end{document}
